\section{Bonus study of parameter estimation with a PINN}
\label{sec:bonus}

\subsection{What this section is for, and how to read its result}
\label{sec:bonus-frame}

The core report solves the \emph{forward} problem: given parameters, predict voltage. This section
reverses it. Given only noisy voltage readings, and treating the recovery variable \(w\) as
unobservable --- which it is in any real recording --- can the model's parameter be recovered?

We answer it twice, with a physics-informed neural network and with a classical least-squares fit,
and we state the expected outcome in advance. Both methods are given the same governing equations
and the same initial state, including the unmeasurable \(w_0\). The least-squares baseline has only
a single scalar to search, with the true model integrated at every trial value, so on this
synthetic problem it is the stronger of the two by construction. It should win, and it does.

The comparison is therefore not a contest but a \textbf{benchmark}: least squares is the empirical
reference point for this experiment, and the quantity of interest is how close a network that must
simultaneously learn two state trajectories comes to it. A PINN landing within a small factor of
that reference is evidence that the implementation is correct. Section~\ref{sec:bonus-value}
returns to the assumptions the baseline relies on and to the settings in which they would not hold.

\subsection{Inverse problem and synthetic data}

The core simulations use known parameters to predict voltage. Here we reverse the problem and estimate \(a\) from voltage readings alone. We fix \(b=0.80\) and \(\varepsilon=0.08\), supply the current and initial state, and treat recovery \(w\) as unobserved during fitting. The data come from the FHN equations with \(a_{\mathrm{true}}=0.70\), so the parameter error can be measured directly.

We generated clean reference trajectories on \([0,40]\) at intervals of \(0.10\), using DOP853 with relative and absolute tolerances of \(10^{-11}\) and \(10^{-13}\). For the clean training set, we kept every fourth voltage sample, giving 101 observations. For the sparse, noisy set, we chose one positive-time sample per block of ten and added Gaussian noise with standard deviation \(0.03\), giving 40 observations. The two sets sample the same training trajectory.

\begin{figure}[H]
  \centering
  \includegraphics[width=0.70\linewidth]{figures/pinn_training_observations.pdf}
  \caption{The two training datasets. Dots are the supplied observations and the gray curve is the clean numerical reference.}
  \label{fig:pinn-data}
\end{figure}

Training used \(I=0.50\) and \((v_0,w_0)=(-1.15,-0.62)\). Validation used \(I=0.65\) and \((-1.00,-0.55)\). The final test used \(I=0.90\) and \((-1.20,-0.70)\). Only training observations entered either fit.

\clearpage

\subsection{PINN and least-squares fitting}

A physics-informed neural network (PINN) fits observed data while penalizing disagreement with the governing equations \cite{raissi2019}. Our network takes time as input and predicts voltage \(\hat v\) and recovery \(\hat w\). It uses normalized time, three pairs of sine and cosine features, and three hidden layers of 32 units with \(\tanh\) activations. The output takes the form \((\hat v,\hat w)=(v_0,w_0)+(1-e^{-t})N_\theta(t)\), which gives the known initial state exactly at \(t=0\).

At 241 fixed physics times, automatic differentiation supplies the predicted slopes. The residuals are the predicted slopes minus the rates from \eqref{eq:voltage} and \eqref{eq:recovery}, evaluated at the predicted state and fitted \(\hat a\). The objective is
\begin{equation}
  \mathcal{L}
  =\operatorname{mean}_{\rm obs}\!\bigl[(\hat v-v_{\rm obs})^2\bigr]
    +\operatorname{mean}_{\rm grid}\!\bigl[r_v^2+4r_w^2\bigr].
  \label{eq:pinn-loss}
\end{equation}
On a Colab CPU in double precision, we used 400 voltage-only steps, 1800 joint Adam steps, and up to 300 L-BFGS iterations. The data seed was 2026 and the network seed was 7.

The least-squares baseline integrates the FHN model for each trial \(a\) and minimizes squared voltage errors at the same observation times. Both fits start at \(a=0.50\) with the same allowed interval \([0.20,1.20]\).

\subsection{Held-out results from the original benchmark}

For validation and testing, we used each fitted \(a\) in the same numerical ODE solver at the new current and initial state. This checks parameter transfer to another trajectory. The voltage error is the root mean square error (RMSE) against the corresponding clean reference. Relative parameter error is \(100|\hat a-a_{\mathrm{true}}|/|a_{\mathrm{true}}|\).

\begin{table}[H]
  \centering
  \caption{Parameter recovery and final-test errors at \(I=0.90\). Entries shown as zero are rounded to the displayed precision.}
  \label{tab:pinn-test}
  \small
  \setlength{\tabcolsep}{5pt}
  \begin{tabular}{llrrrr}
    \toprule
    Training data & Method & \(\hat a\) & Rel.\ error & Voltage RMSE & Recovery RMSE\\
    \midrule
    Clean & PINN          & 0.699975 & \(0.0036\%\) & 0.000096 & 0.000022\\
    Clean & Least squares & 0.700000 & \(0.00\%\)   & 0.000000 & 0.000000\\
    Sparse, noisy & PINN          & 0.704093 & \(0.58\%\) & 0.015570 & 0.003554\\
    Sparse, noisy & Least squares & 0.701926 & \(0.28\%\) & 0.007343 & 0.001677\\
    \bottomrule
  \end{tabular}
\end{table}

\begin{figure}[H]
  \centering
  \includegraphics[width=0.91\linewidth]{figures/pinn_heldout_comparison.pdf}
  \caption{Final-test voltages and absolute errors for the sparse, noisy fits. The curves are regenerated from the fitted parameters rounded to six decimal places.}
  \label{fig:pinn-test}
\end{figure}

Both methods matched the single reference spike in the final test. On the clean validation trajectory, the fits from noisy training readings gave voltage RMSEs of \(0.011649\) for the PINN and \(0.005491\) for least squares.


\clearpage

\subsection{Controlled comparison of sampling and noise}
\label{sec:controlled}

The original benchmark changed the number of readings and the amount of noise together. We now vary these two factors separately on the same training trajectory. The dense grid uses 100 positive-time samples from \(0.4\) through \(40\), spaced by \(0.4\) model time units. The sparse grid keeps the second and fifth reading from each block of five dense readings, for a total of 40 samples. The known initial state at \(t=0\) is shared by all fits.

For each noise seed \(2026,\dots,2030\), we draw independent Gaussian voltage errors with mean zero and standard deviation \(0.03\) on the dense grid. The sparse dataset uses the same perturbed readings at its selected times. Each density also has a clean control at the same times. This yields two clean datasets and ten noisy datasets. Related work has studied the sensitivity of FHN parameter estimation to observation noise with other neural architectures \cite{rudi2022}.

Both methods use the same readings, bounds, and known initial state. The PINN architecture, loss, collocation grid, optimizer budget, and initialization are held fixed. Each fitted \(\hat a\) is scored by relative parameter error and by voltage RMSE on the final trajectory at \(I=0.90\). That trajectory is used only for scoring.

\paragraph{How the comparisons isolate each factor.}
For noise, subtract the clean error from the noisy error at the same sample count. For sampling, subtract the dense error from the sparse error using the same noise seed. Positive differences indicate an increase in error. We report medians and full ranges across the five noise realizations. The ranges describe the tested realizations.

\begin{table}[H]
  \centering
  \caption{Controlled experiment results on the final test at \(I=0.90\). Noisy conditions show the median and full range in brackets. The two clean datasets contribute one fit per method.}
  \label{tab:controlled-both}
  \small
  \input{pinn_controlled_results}
\end{table}

Both methods recovered \(a\) closely from clean readings. For the noisy datasets, least squares gave smaller parameter errors and lower held-out voltage RMSEs in all ten matched comparisons. At 100 readings, the median RMSE was \(0.002744\) for least squares and \(0.018354\) for the PINN. At 40 readings, the corresponding medians were \(0.002029\) and \(0.022869\). The least-squares fit solves the exact data-generating equations at each trial value of \(a\). The PINN estimates \(a\) while also learning two state curves and balancing its voltage and equation losses. This benchmark compares those two procedures under a known model and fixed training settings.

\clearpage

\subsection{Separate effects of noise and sample count}

\begin{figure}[H]
  \centering
  \includegraphics[width=\linewidth]{figures/pinn_controlled_comparison.pdf}
  \caption{Parameter error and held-out voltage RMSE by observation condition. Each marker is one fit. The horizontal bars show group medians. All held-out trajectories use the same ODE solver with the fitted parameter.}
  \label{fig:controlled-both}
\end{figure}

\begin{table}[H]
  \centering
  \caption{Paired changes in held-out voltage RMSE. Each noisy contrast summarizes five differences as median [minimum, maximum]. Positive values mean error increased.}
  \label{tab:controlled-paired}
  \small
  \begin{tabular}{@{}lll@{}}
    \toprule
    Change & Least squares & PINN\\
    \midrule
    Add noise at 100 readings & \(+0.002744\) [0.001384, 0.006811] & \(+0.018283\) [0.007749, 0.022230]\\
    Add noise at 40 readings & \(+0.002029\) [0.000996, 0.007552] & \(+0.022628\) [0.004908, 0.030079]\\
    Reduce noisy readings to 40 & \(+0.000645\) [$-0.005139$, 0.004808] & \(+0.008018\) [$-0.008281$, 0.012244]\\
    \bottomrule
  \end{tabular}
\end{table}

At a fixed set of observation times, adding noise raised the test error in every trial for both methods. With noisy readings, reducing the sample count increased the error in three of five least-squares pairs and four of five PINN pairs. The effect of the sample count therefore depended on which readings were retained and on their noise values. For example, least squares had a lower group median with 40 noisy readings even though its median \emph{paired change} was positive. The paired change is calculated by subtracting errors seed by seed before taking the median.

These comparisons summarize five realizations of one noise level. They show performance for the specified FHN model, observation times, and fixed initialization. Other noise levels, observation patterns, or initializations would require separate trials.

\clearpage

\subsection{What this section contributes, and where it stops}
\label{sec:bonus-value}

\begin{table}[H]
  \centering
  \caption{What the bonus study delivers, and where each item is evidenced.}
  \label{tab:bonus-map}
  \small
  \begin{tabular}{@{}p{0.27\linewidth}p{0.66\linewidth}@{}}
    \toprule
    Element & Where, and what was actually done\\
    \midrule
    Advanced numerics &
      Automatic differentiation supplies exact derivatives of the network at 241 collocation
      times; a three-stage schedule of 400 data-only steps, 1800 joint Adam steps and up to 300
      L-BFGS iterations with strong-Wolfe line search. The hard-constraint ansatz
      \((v_0,w_0)+(1-e^{-t})N_\theta(t)\) satisfies the initial condition exactly rather than by
      penalty, removing a loss term and its weight.\\[0.35em]
    Machine learning &
      A physics-informed network estimates the unknown parameter as a trainable weight while
      learning the state curves, from voltage observations alone with \(w\) hidden.\\[0.35em]
    Validation &
      Each fitted parameter is scored on \emph{held-out trajectories at currents never seen in
      training}: validation at \(I=0.65\) and a final test at \(I=0.90\), both from different
      initial states, re-integrated with the ODE solver. Both methods reproduce the correct spike
      count on the held-out run.\\[0.35em]
    Controlled experiment &
      The first benchmark confounded noise with sample count. Section~\ref{sec:controlled}
      replaces it with a paired design: two clean datasets and ten noisy ones (two densities
      \(\times\) five noise seeds), sparse readings copied from the dense ones at shared times,
      one fixed network initialisation. Each of the twelve datasets is fitted by both methods,
      giving twenty-four fits, reported in full.\\[0.35em]
    Interpretation &
      The negative result is explained rather than reported: the baseline is identified as the
      empirical reference for this synthetic experiment, so the PINN's distance from it is
      the quantity of interest.\\
    \bottomrule
  \end{tabular}
\end{table}

\paragraph{Proposed follow-up experiments.}
The baseline's advantage here rests on three assumptions: the model is correct, the initial state
is known, and the unknowns are finitely many scalars. Each can be relaxed, and doing so is the
natural continuation of this study. We have not run these experiments; we set them out as the
tests that would determine whether the ordering above survives.

\begin{itemize}[leftmargin=1.4em]
  \item \textbf{Unknown forcing.} If \(I(t)\) is an unknown \emph{function} rather than a
        constant, a second network head could estimate it pointwise, whereas the baseline would
        first have to assume a parametric form. Whether that changes the ordering is untested.
  \item \textbf{Unknown initial state.} Here both methods were given \(w_0\). In a recording it
        is unobservable, and the baseline would have to search over it too. The experiment is to
        withhold \(w_0\) from both and re-measure the gap.
  \item \textbf{Model misspecification.} With a deliberately wrong right-hand side, the question
        is whether the network's physics residual grows in a way that diagnoses the error while
        the baseline absorbs it into the fitted parameter.
  \item \textbf{Expensive forward solves.} For a spatially extended version of these equations,
        the cost of re-integrating a PDE inside the optimiser at every trial parameter would need
        to be compared against the cost of the collocation formulation.
\end{itemize}

\paragraph{Limitations, stated plainly.}
All observations are synthetic. That is deliberate --- the true
parameter is known only for synthetic data and therefore to measure error rather than merely report a fit --- but it means the
study demonstrates method behaviour, not biological fidelity. It uses one noise level, five noise
realisations, one network initialisation and one unknown parameter. Fitting a published
patch-clamp recording, estimating several parameters jointly, and sweeping the noise level to look
for a crossover are the natural continuations, and none of them is addressed here.
